\subsection{Normed and Banach Spaces}

% **Role:** infinite-dimensional linear algebra powered by analysis.
%
% Normed spaces, Banach spaces, bounded linear maps, operator norms, examples from sequence/function spaces.

Normed spaces make vector spaces metric; Banach spaces add completeness. This is the setting for infinite series, fixed points, bounded linear maps, and operator norms.

\begin{definition} \label{def:normed-and-banach-spaces} \hlt{Normed and Banach Spaces}\\
Let $\mathcal{X}$ be a vector space over $\mathbb{F}\in\{\R,\C\}$.
\begin{enumerate}[label=(\roman*)]
\setlength{\itemsep}{0pt}
\item \hlt{Norm}: $\norm{\cdot}:\mathcal{X}\to[0,\infty)$; $\norm{x}=0\iff x=0$, $\norm{\lambda x}=\abs{\lambda}\norm{x}$; $\norm{x+y}\leq\norm{x}+\norm{y}$ for $x,y\in\mathcal{X}$, $\lambda\in\mathbb{F}$.
\item \hlt{Normed space}: $(\mathcal{X},\norm{\cdot})$ carries the metric $d(x,y)\defeq\norm{x-y}$.
\item \hlt{Banach space}: $(\mathcal{X},\norm{\cdot})$ is complete in the norm metric.
\end{enumerate}
\end{definition}

\begin{definition} \label{def:basic-normed-space-geometry} \hlt{Basic Normed-Space Geometry}\\
Let $\mathcal{X}$ be a normed space.
\begin{enumerate}[label=(\roman*)]
\setlength{\itemsep}{0pt}
\item \hlt{Open unit ball}: $B_{\mathcal{X}}\defeq\{x\in\mathcal{X}:\norm{x}<1\}$.
\item \hlt{Closed unit ball}: $\overline{B}_{\mathcal{X}}\defeq\{x\in\mathcal{X}:\norm{x}\leq1\}$.
\item \hlt{Unit vector}: $x\in\mathcal{X}$ satisfies $\norm{x}=1$.
\item \hlt{Minkowski sum}: $A+B\defeq\{a+b:a\in A,\ b\in B\}$ for nonempty $A,B\subseteq\mathcal{X}$.
\item \hlt{Scalar multiple}: $\lambda A\defeq\{\lambda a:a\in A\}$ for nonempty $A\subseteq\mathcal{X}$ and $\lambda\in\mathbb{F}$.
\item \hlt{Closed linear span}: $\operatorname{cspan}(S)\defeq\overline{\operatorname{span}(S)}$ for $S\subseteq\mathcal{X}$.
\item \hlt{Equivalent norms}: $\norm{\cdot}_1\sim\norm{\cdot}_2$ iff $\exists c_1,c_2>0$ such that $\forall x\in\mathcal{X}$, $c_1\norm{x}_1\leq\norm{x}_2\leq c_2\norm{x}_1$.
\end{enumerate}
\end{definition}

\begin{proposition} \label{prop:linear-operations-closed-spans-and-equivalent-norms} \hlt{Linear Operations, Closed Spans, and Equivalent Norms}\\
Let $\mathcal{X}$ be a normed space.
\begin{enumerate}[label=(\roman*)]
\setlength{\itemsep}{0pt}
\item The norm map $x\mapsto\norm{x}$, addition $\mathcal{X}\times\mathcal{X}\to\mathcal{X}$, and scalar multiplication $\mathbb{F}\times\mathcal{X}\to\mathcal{X}$ are continuous.
\item If $M\leq\mathcal{X}$, then $\overline{M}\leq\mathcal{X}$; for $S\subseteq\mathcal{X}$, $\operatorname{cspan}(S)$ is the least closed subspace containing $S$.
\item Equivalent norms have the same convergent and Cauchy sequences, topology, and compact sets.
\end{enumerate}
\begin{proof}
\begin{enumerate}[label=(\roman*)]
\setlength{\itemsep}{0pt}
\item $\abs{\norm{x}-\norm{y}}\leq\norm{x-y}$, so $x\mapsto\norm{x}$ is $1$-Lipschitz.\\
Also $\norm{(x+h)+(y+k)-(x+y)}\leq\norm h+\norm k$, so addition is continuous.\\
If $\lambda_n\to\lambda$ and $x_n\to x$, then $\norm{\lambda_nx_n-\lambda x}\leq\abs{\lambda_n}\norm{x_n-x}+\abs{\lambda_n-\lambda}\norm x\to0$, since $(\lambda_n)$ is bounded.
\item Let $x,y\in\overline M$ and choose $x_n,y_n\in M$ with $x_n\to x$, $y_n\to y$.\\
For $\alpha,\beta\in\mathbb F$, $\alpha x_n+\beta y_n\in M$ and $\alpha x_n+\beta y_n\to\alpha x+\beta y$ by (i); hence $\overline M\leq\mathcal X$.\\
If $C$ is closed, $C\leq\mathcal X$, and $S\subseteq C$, then $\operatorname{span}(S)\subseteq C$, so $\operatorname{cspan}(S)=\overline{\operatorname{span}(S)}\subseteq C$.
\item If $c\norm{z}_1\leq\norm{z}_2\leq C\norm{z}_1$ for all $z\in\mathcal X$, then $B_1(x,\frac{\varepsilon}{C})\subseteq B_2(x,\varepsilon)\subseteq B_1(x,\frac{\varepsilon}{c})$. Thus the topologies agree, the identity maps are bi-Lipschitz, convergent and Cauchy sequences agree, and compact sets agree.
\end{enumerate}
\hlt{Structural Commentary}: Norms turn the algebraic operations into continuous operations.\\
\hlt{Key Insight}: Closing a subspace preserves linearity because linear combinations commute with limits.
\end{proof}
\end{proposition}

\begin{definition} \label{def:bounded-linear-maps-and-operator-norm} \hlt{Bounded Linear Maps and Operator Norm}\\
Let $\mathcal{X},\mathcal{Y}$ be normed spaces and let $T:\mathcal{X}\to\mathcal{Y}$ be linear.
\begin{enumerate}[label=(\roman*)]
\setlength{\itemsep}{0pt}
\item \hlt{Bounded}: $\exists C\geq0$ such that $\forall x\in\mathcal{X}$, $\norm{Tx}_{\mathcal{Y}}\leq C\norm{x}_{\mathcal{X}}$.
\item \hlt{Operator norm}: $\normop{T}\defeq\sup\limits_{\norm{x}_{\mathcal{X}}\leq1}\norm{Tx}_{\mathcal{Y}}=\sup\limits_{x\neq0}\frac{\norm{Tx}_{\mathcal{Y}}}{\norm{x}_{\mathcal{X}}}$.
\item \hlt{Operator space}: $\LL(\mathcal{X},\mathcal{Y})\defeq\{T:\mathcal{X}\to\mathcal{Y}:T\text{ is bounded and linear}\}$ and $\LL(\mathcal{X})\defeq\LL(\mathcal{X},\mathcal{X})$.
\end{enumerate}
\end{definition}

\begin{theorem} \label{thm:bounded-linear-map-continuity-equivalences} \hlt{Bounded Equals Continuous for Linear Maps}\\
For a linear map $T:\mathcal{X}\to\mathcal{Y}$ between normed spaces, the following are equivalent:
\begin{enumerate}[label=(\roman*)]
\setlength{\itemsep}{0pt}
\item $T$ is bounded.
\item $T$ is continuous at $0$.
\item $T$ is continuous everywhere.
\end{enumerate}
\begin{proof}
\begin{enumerate}[label=(\roman*)]
\setlength{\itemsep}{0pt}
\item[$(i)$]$\Rightarrow (iii)$ If $T$ is bounded, then $\norm{Tx-Ty}\leq\normop{T}\norm{x-y}$ for $x,y\in\mathcal X$; hence $T$ is Lipschitz.
\item[$(iii)$]$\Rightarrow (ii)$ Since $0\in\mathcal X$, continuity on $\mathcal X$ implies continuity at $0$.
\item[$(ii)$]$\Rightarrow (i)$ Choose $\delta>0$ with $\norm z<\delta\Rightarrow\norm{Tz}<1$. For $x\ne0$, set $z\defeq\frac{\delta x}{2\norm{x}}$.\\
Then $\norm z<\delta$ and $\norm{Tx}=\frac{2\norm{x}}{\delta}\norm{Tz}\leq\frac{2}{\delta}\norm{x}$. Thus $\normop T\leq\frac{2}{\delta}$.
\end{enumerate}
\hlt{Structural Commentary}: Continuity at $0$ rescales to a global operator bound.\\
\hlt{Key Insight}: Linear continuity is boundedness.
\end{proof}
\end{theorem}

\begin{theorem} \label{thm:operator-space-banach} \hlt{Bounded Operators Form a Banach Space}\\
Let $\mathcal{X},\mathcal{Y}$ be normed spaces.
\begin{enumerate}[label=(\roman*)]
\setlength{\itemsep}{0pt}
\item $\LL(\mathcal{X},\mathcal{Y})$ is normed by $\normop{\cdot}$ under pointwise operations.
\item If $\mathcal{Y}$ is Banach, then $\LL(\mathcal{X},\mathcal{Y})$ is Banach.
\end{enumerate}
\begin{proof}
\begin{enumerate}[label=(\roman*)]
\setlength{\itemsep}{0pt}
\item If $S,T\in\LL(\mathcal X,\mathcal Y)$ and $\alpha,\beta\in\mathbb F$, then $\norm{(\alpha S+\beta T)x}\leq(\abs\alpha\normop S+\abs\beta\normop T)\norm x$, so $\alpha S+\beta T\in\LL(\mathcal X,\mathcal Y)$.\\
Also $\normop{\alpha T}=\abs\alpha\normop T$, $\normop{S+T}\leq\normop S+\normop T$, and $\normop T=0\Rightarrow Tx=0$ for all $x\in\mathcal X$.
\item Assume $\mathcal Y$ is Banach and let $(T_n)$ be Cauchy in $\normop{\cdot}$.\\
For each $x\in\mathcal X$, $\norm{T_nx-T_mx}\leq\normop{T_n-T_m}\norm x$, so define $Tx\defeq\lim\limits_nT_nx$.\\
Then $T(\alpha x+\beta y)=\alpha Tx+\beta Ty$. Given $\varepsilon>0$, choose $N$ with $\normop{T_n-T_m}<\varepsilon$ for $m,n\geq N$.\\
Letting $m\to\infty$ gives $\norm{(T_n-T)x}\leq\varepsilon\norm x$, so $\normop{T_n-T}\leq\varepsilon$.\\
Hence $T_N-T\in\LL(\mathcal X,\mathcal Y)$, $T\in\LL(\mathcal X,\mathcal Y)$, and $T_n\to T$.
\end{enumerate}
\hlt{Structural Commentary}: Pointwise limits are linear; operator-norm Cauchy control is uniform in $x$.\\
\hlt{Key Insight}: Completeness of $\mathcal Y$ lifts to $\LL(\mathcal X,\mathcal Y)$.
\end{proof}
\end{theorem}

\begin{corollary} \label{cor:neumann-series-banach-operator} \hlt{Neumann Series}\\
Let $\mathcal{X}$ be Banach and let $T\in\LL(\mathcal{X})$ satisfy $\normop{T}<1$.\\
Then $I-T$ has bounded inverse, with Neumann series converging in $\LL(\mathcal{X})$:
$(I-T)^{-1}=\sum\limits_{n=0}^{\infty}T^n$.
\begin{proof} Let $S_N\defeq\sum\limits_{n=0}^{N}T^n$. For $M>N$, $\normop{S_M-S_N}\leq\sum\limits_{n=N+1}^{M}\normop T^n\leq\frac{\normop T^{N+1}}{1-\normop T}\to0$, so $S_N\to S\in\LL(\mathcal X)$.\\
Since $(I-T)S_N=I-T^{N+1}=S_N(I-T)$ and $\normop{T^{N+1}}\leq\normop T^{N+1}\to0$.\\
Passing to the limit gives $(I-T)S=S(I-T)=I$.\\
\hlt{Structural Commentary}: Operator-norm convergence turns the geometric identity into an inverse.\\
\hlt{Key Insight}: If $\normop T<1$, then $\sum T^n=(I-T)^{-1}$.
\end{proof}
\end{corollary}

\begin{proposition} \label{prop:quotient-norm-and-map} \hlt{Quotient Norm and Quotient Map}\\
Let $M\leq\mathcal X$ be closed, let $\mathcal X/M\defeq\{x+M:x\in\mathcal X\}$, and let $q:\mathcal X\to\mathcal X/M$ be $q(x)\defeq x+M$.\\
Set $\norm{x+M}_{\mathcal X/M}\defeq\inf\limits_{m\in M}\norm{x-m}$.
\begin{enumerate}[label=(\roman*)]
\setlength{\itemsep}{0pt}
\item $\norm{\cdot}_{\mathcal{X}/M}$ is a norm on $\mathcal{X}/M$.
\item $q$ is continuous with $\normop{q}\leq1$.
\item $\forall x\in\mathcal{X}$ and $r>0$, $q(B_r(x))=B_r(qx)$.
\item $q$ is open.
\end{enumerate}
\begin{proof}
\begin{enumerate}[label=(\roman*)]
\setlength{\itemsep}{0pt}
\item If $x'=x+m_0$ with $m_0\in M$, then $\inf\limits_{m\in M}\norm{x'-m}=\inf\limits_{m\in M}\norm{x-m}$, so the formula is well defined.\\
Also $\norm{x+M}=0\iff x\in\overline M=M$. For $\alpha\ne0$,\\
$\norm{\alpha(x+M)}=\inf\limits_{m\in M}\norm{\alpha x-m}=\abs\alpha\inf\limits_{m\in M}\norm{x-\alpha^{-1}m}=\abs\alpha\norm{x+M}$; $\alpha=0$ is trivial.\\
For $m,n\in M$, $\norm{(x+y)-(m+n)}\leq\norm{x-m}+\norm{y-n}$; take infima.
\item Since $\norm{q(x)}\leq\norm{x-0}=\norm{x}$, $q\in\LL(\mathcal X,\mathcal X/M)$ and $\normop q\leq1$.
\item If $z\in B_r(x)$, then $\norm{qz-qx}\leq\norm{z-x}<r$, so $q(B_r(x))\subseteq B_r(qx)$. Conversely, if $\eta=qz$ and $\norm{\eta-qx}<r$, choose $m\in M$ with $\norm{z-x-m}<r$; then $z-m\in B_r(x)$ and $q(z-m)=\eta$.
\item If $U\subseteq\mathcal X$ is open and $y\in q(U)$, choose $x\in U$ with $qx=y$ and $r>0$ with $B_r(x)\subseteq U$.\\
By (iii), $B_r(y)=q(B_r(x))\subseteq q(U)$; hence $q(U)$ is open.
\end{enumerate}
\hlt{Structural Commentary}: Quotient norm is distance to the closed subspace $M$.\\
\hlt{Key Insight}: Closedness makes zero distance equivalent to membership.
\end{proof}
\end{proposition}

\begin{theorem} \label{thm:banach-quotient-space} \hlt{Banach Quotients}\\
If $M$ is a closed linear subspace of a Banach space $\mathcal{X}$, then $\mathcal{X}/M$ is Banach in the quotient norm.
\begin{proof}
Let $(\xi_n)$ be Cauchy in $\mathcal X/M$. Choose $n_1<n_2<\cdots$ with $\norm{\xi_{n_{k+1}}-\xi_{n_k}}<2^{-k}$ and choose $u_k\in\mathcal X$ with $q(u_k)=\xi_{n_{k+1}}-\xi_{n_k}$, $\norm{u_k}<2^{-k}$. Since $\sum\limits_k\norm{u_k}<+\infty$ and $\mathcal X$ is Banach, $u\defeq\sum\limits_{k\geq1}u_k$ exists.\\
For $K\geq1$, $\xi_{n_K}=\xi_{n_1}+q(\sum\limits_{k<K}u_k)\to\xi_{n_1}+q(u)$.\\
Since $(\xi_n)$ is Cauchy and has a convergent subsequence, $\xi_n\to\xi_{n_1}+q(u)$.\\
\hlt{Structural Commentary}: Rapid quotient increments lift to an absolutely convergent series.\\
\hlt{Key Insight}: A Cauchy sequence with a convergent subsequence converges.
\end{proof}
\end{theorem}

\begin{proposition} \label{prop:kernel-quotient-factor-operator-norm} \hlt{Operator Factor through the Kernel Quotient}\\
Let $T\in\LL(\mathcal{X},\mathcal{Y})$, set $M\defeq\ker T$, and let $\widetilde{T}:\mathcal{X}/M\to\mathcal{Y}$ be induced by Theorem \ref{thm:quotient-by-kernel-is-range}.
\begin{enumerate}[label=(\roman*)]
\setlength{\itemsep}{0pt}
\item $\widetilde{T}$ is bounded.
\item $\normop{\widetilde{T}}=\normop{T}$.
\end{enumerate}
\begin{proof}
Since $q(x)=q(y)\iff x-y\in M=\ker T$, the formula $\widetilde T(qx)\defeq Tx$ is well defined and $T=\widetilde Tq$.\\
For $m\in M$, $\norm{\widetilde T(qx)}=\norm{Tx}=\norm{T(x-m)}\leq\normop T\norm{x-m}$; taking $\inf\limits_{m\in M}$ gives $\norm{\widetilde T(qx)}\leq\normop T\norm{qx}$.
\begin{enumerate}[label=(\roman*)]
\setlength{\itemsep}{0pt}
\item The estimate gives $\widetilde T\in\LL(\mathcal X/M,\mathcal Y)$.
\item The estimate gives $\normop{\widetilde T}\leq\normop T$. Since $T=\widetilde Tq$ and $\normop q\leq1$, $\normop T\leq\normop{\widetilde T}\normop q\leq\normop{\widetilde T}$.
\end{enumerate}
\hlt{Structural Commentary}: Collapsing $\ker T$ makes $T$ descend without changing its norm.\\
\hlt{Key Insight}: The quotient infimum is taken over exactly the invisible directions of $T$.
\end{proof}
\end{proposition}

\begin{example} \label{ex:standard-banach-space-examples} \hlt{Standard Banach Space Examples}\\
Standard Banach spaces include the following canonical examples.
\begin{enumerate}[label=(\roman*)]
\setlength{\itemsep}{0pt}
\item If $S$ is a set and $B(S)\defeq\{f:S\to\mathbb F:\norm f_\infty<+\infty\}$, then $(B(S),\norm{\cdot}_\infty)$ is Banach.
\item $\ell^\infty\defeq\{x=(x_n):\sup\limits_n\abs{x_n}<+\infty\}$ is Banach under $\norm{x}_\infty\defeq\sup\limits_n\abs{x_n}$.
\item For $1\leq p<\infty$, $\ell^p\defeq\{x=(x_n):\sum\limits_n\abs{x_n}^p<+\infty\}$ is Banach under $\norm{x}_p\defeq(\sum\limits_n\abs{x_n}^p)^{1/p}$.
\item For a measure space $(X,\mathcal M,\mu)$ and $1\leq p\leq\infty$, $L^p(\mu)$ is Banach.
\item $L^2(\mu)$ is Hilbert with $\langle f,g\rangle\defeq\int_X f\overline g\,d\mu$.
\end{enumerate}
\end{example}

\begin{theorem} \label{thm:finite-dimensional-normed-spaces} \hlt{Finite-Dimensional Normed Spaces}\\
Let $\mathcal{X}$ be a finite-dimensional normed space.
\begin{enumerate}[label=(\roman*)]
\setlength{\itemsep}{0pt}
\item $\mathcal{X}$ is Banach.
\item Every subspace of $\mathcal{X}$ is closed.
\item $\overline{B}_{\mathcal{X}}$ is compact.
\end{enumerate}
\begin{proof}
Choose a basis $e_1,\ldots,e_n$ and set $A:\mathbb F^n\to\mathcal X$, $Aa\defeq\sum\limits_{i=1}^{n}a_ie_i$.\\
Then $A$ is linear bijective and $\norm{Aa}\leq\sum\limits_i\abs{a_i}\norm{e_i}\leq C\norm a_2$.\\
Since $a\mapsto\norm{Aa}$ is continuous, positive on $\{a:\norm a_2=1\}$, $\exists c>0$ with $c\norm a_2\leq\norm{Aa}\leq C\norm a_2$ $\forall a\in\mathbb F^n$.
\begin{enumerate}[label=(\roman*)]
\setlength{\itemsep}{0pt}
\item If $(x_k)$ is Cauchy in $\mathcal X$, then $a_k\defeq A^{-1}x_k$ is Cauchy in $\mathbb F^n$ by $c\norm{a_k-a_\ell}_2\leq\norm{x_k-x_\ell}$.\\
Since $\mathbb F^n$ is complete, $a_k\to a$, hence $x_k\to Aa$.
\item If $M\leq\mathcal X$, then $M$ is finite-dimensional, hence Banach by (i). If $x_k\in M$ and $x_k\to x$ in $\mathcal X$, then $(x_k)$ converges in $M$ to some $y\in M$; uniqueness of limits gives $x=y\in M$. Thus $M$ is closed.
\item Since $\overline{B}_{\mathcal X}\subseteq A(\overline{B}_2(0,\frac{1}{c}))$ and $A^{-1}(\overline{B}_{\mathcal X})$ is closed in the compact Euclidean ball $\overline{B}_2(0,\frac{1}{c})$, $A^{-1}(\overline{B}_{\mathcal X})$ is compact; continuity of $A$ gives compactness of $\overline{B}_{\mathcal X}$.
\end{enumerate}
\hlt{Structural Commentary}: The coordinate map is a bi-Lipschitz linear isomorphism.\\
\hlt{Key Insight}: Finite-dimensional normed spaces inherit completeness and compactness from $\mathbb F^n$.
\end{proof}
\end{theorem}

\begin{theorem} \label{thm:riesz-lemma-unit-ball-compactness} \hlt{Riesz Lemma}\\
Let $M$ be a closed proper linear subspace of a normed space $\mathcal{X}$.\\
Then $\forall\varepsilon\in(0,1)\ \exists u\in\mathcal{X}$ such that $\norm{u}=1$ and $\forall m\in M$, $\norm{u-m}>1-\varepsilon$.
\begin{proof}
Choose $x\in\mathcal X\setminus M$ and set $d\defeq\inf\limits_{m\in M}\norm{x-m}$. Since $M$ is closed, $d>0$.\\
Choose $m_0\in M$ with $d\leq\norm{x-m_0}<\frac{d}{1-\varepsilon}$ and set $u\defeq\frac{x-m_0}{\norm{x-m_0}}$.\\
Then $\norm u=1$. For $m\in M$, $m_0+\norm{x-m_0}m\in M$, so $\norm{u-m}=\frac{\norm{x-(m_0+\norm{x-m_0}m)}}{\norm{x-m_0}}\geq\frac{d}{\norm{x-m_0}}>1-\varepsilon$.\\
\hlt{Structural Commentary}: A positive distance to $M$ survives normalization.\\
\hlt{Key Insight}: Normalize an almost-minimizer of $\operatorname{dist}(x,M)$.
\end{proof}
\end{theorem}

\begin{corollary} \label{cor:compact-unit-ball-finite-dimensional} \hlt{Compact Unit Balls Characterize Finite Dimension}\\
For a normed space $\mathcal{X}$, $\overline{B}_{\mathcal{X}}$ is compact iff $\dim\mathcal{X}<\infty$.
\begin{proof}
\begin{enumerate}[label=\roman*.]
\setlength{\itemsep}{0pt}
\item[$(\Leftarrow)$] If $\dim\mathcal X<\infty$, then $\overline B_{\mathcal X}$ is compact by Theorem \ref{thm:finite-dimensional-normed-spaces}(iii).
\item[$(\Rightarrow)$] Suppose $\dim\mathcal X=\infty$. Choose $x_1\in\mathcal X$ with $\norm{x_1}=1$. Given $x_1,\ldots,x_n$, set $M_n\defeq\operatorname{span}\{x_1,\ldots,x_n\}$.\\
Then $M_n$ is finite-dimensional, closed, and proper, so Theorem \ref{thm:riesz-lemma-unit-ball-compactness} with $\varepsilon=\frac{1}{2}$ gives $x_{n+1}\in\mathcal X$ with $\norm{x_{n+1}}=1$ and $\operatorname{dist}(x_{n+1},M_n)>\frac{1}{2}$. Since $x_i\in M_n$ for $i\leq n$, $\norm{x_{n+1}-x_i}>\frac{1}{2}$.\\
Thus $(x_n)\subseteq\overline B_{\mathcal X}$ has no Cauchy subsequence, so $\overline B_{\mathcal X}$ is not compact.
\end{enumerate}
\hlt{Structural Commentary}: Riesz iteration builds a separated sequence in $\overline B_{\mathcal X}$.\\
\hlt{Key Insight}: Compact unit ball forces finite dimension.
\end{proof}
\end{corollary}

\begin{figure}[H]
  \centering
  \begin{tikzpicture}[
  >=Stealth,
  every node/.style={font=\scriptsize},
  ball/.style={draw=blue!55!black, fill=blue!5, line width=0.65pt},
  subspace/.style={line width=0.65pt, black!55},
  vector/.style={->, line width=0.65pt, black!65},
  point/.style={circle, fill=orange!85!black, inner sep=1.35pt},
  sep/.style={draw=orange!80!black, dashed, line width=0.5pt},
  arrow/.style={->, line width=0.55pt, black!65, shorten <=2pt, shorten >=2pt},
  note/.style={fill=white, inner sep=1.1pt, text=black!72, align=center}
]
  \begin{scope}[shift={(-3.45,0)}]
    \draw[ball] (0,0) circle (1.0);
    \draw[subspace] (-1.40,-0.40) -- (1.40,0.40);
    \node[note, anchor=east] at (-1.28,-0.58) {$M$};

    \coordinate (u) at (0.38,0.93);
    \coordinate (m) at (0.50,0.14);
    \draw[vector] (0,0) -- (u);
    \node[point] at (u) {};
    \node[note, anchor=south west] at (0.52,1.02) {$u$};
    \draw[sep] (u) -- (m);
    \node[note, anchor=west] at (1.12,0.64)
      {$\operatorname{dist}(u,M)>1-\varepsilon$};
  \end{scope}

  \draw[arrow] (-1.58,0) -- (1.58,0);

  \begin{scope}[shift={(3.45,0)}]
    \draw[ball] (0,0) circle (1.0);
    \foreach \ang/\lab in {18/1,82/2,148/3,216/4,292/5} {
      \coordinate (p\lab) at ({0.91*cos(\ang)},{0.91*sin(\ang)});
      \node[point] at (p\lab) {};
      \node[font=\tiny, black!70] at ({1.15*cos(\ang)},{1.15*sin(\ang)}) {$x_{\lab}$};
    }
    \draw[sep] (p1) -- (p2);
    \draw[sep] (p2) -- (p3);
    \draw[sep] (p3) -- (p4);
    \draw[sep] (p4) -- (p5);
    \node[note] at (0,1.24) {$\norm{x_i-x_j}>\frac{1}{2}$};
  \end{scope}
\end{tikzpicture}

  \caption{Riesz iteration: $M_n=\operatorname{span}\{x_1,\ldots,x_n\}$, $\operatorname{dist}(x_{n+1},M_n)>\frac{1}{2}\Rightarrow \norm{x_i-x_j}>\frac{1}{2}$ $(i\neq j)$; hence no Cauchy subsequence in $\overline{B}_{\mathcal X}$.}
  \label{fig:riesz-lemma-separated-unit-ball}
\end{figure}
